% HD-F1-T27. GPT-6 Astra Pro; /root/transport; 22 November 2026 UTC.
\section[HD-F1-T27: all third-tail Gaussian Specht sources are excluded]{HD-F1-T27: the original permutation quotient excludes all third-tail Gaussian Specht families}
\label{sec:f1-gaussian-three-tail-exclusion}

In the residual Gaussian source of T24--T26, none of
\[
 (n-3,3),\quad(n-3,2,1),\quad(n-3,1,1,1)
 \quad\hbox{or their conjugate partitions}
 \tag{F1T27.1}
\]
occurs. Thus its partition has at least four boxes outside
both its first row and its first column. A separate elementary
dimension argument below also gives the explicit retained bound
\[
             m\ge n(n-1)(n-5)/6\ge1309.
 \tag{F1T27.2}
\]
Here $n\ge22$ is the already proved T25 bound. The exclusion
uses the original branch tuple, its scalar character, and its
actual degree-$n$ quotient; none is replaced.

\paragraph{1. Exact character formulas and uniform affine bounds.}
Write the original coefficient as
$\rho(g)=c(g)R_\lambda(\sigma(g))$, as constructed in T24,
with $c(g)\in\mu_4$ and $\sigma(G)=A_n$ or $S_n$.
For a permutation $s$, let $f,k_2,k_3$ be its numbers of
cycles of lengths one, two and three. The exact character
formulas on these actual rational Specht modules are
\[
\begin{array}{c|c|c}
\lambda& m=\chi(1)&\chi(s)\\ \hline
(n-3,3)&n(n-1)(n-5)/6&
\binom f3-\binom f2+(f-1)k_2+k_3\\
(n-3,1^3)&(n-1)(n-2)(n-3)/6&
\binom{f-1}3-(f-1)k_2+k_3\\
(n-3,2,1)&n(n-2)(n-4)/3&
f(f-2)(f-4)/3-k_3 .
\end{array}
 \tag{F1T27.3}
\]
Binomial symbols here mean falling-factorial polynomials,
also when $f=0$. These are Guralnick--Tiep, Lemma3.3,
whose complete supplied proof was read. In particular this
is a trace formula for the existing representation, not
an assertion that a virtual character is an actual replacement.

For every nonidentity $s$ the following bound holds:
\[
 |\chi(s)|\le Af+B,\qquad
\begin{array}{c|c|c|c}
\lambda&A&B&\beta:=m-An-B\\ \hline
(n-3,3)&(n-4)(n-5)/6&n/2&n(n-6)/2\\
(n-3,1^3)&(n-4)(n-5)/6&n/2&(n^2-4n-2)/2\\
(n-3,2,1)&(n-4)(n-6)/3&n/3&n(4n-17)/3 .
\end{array}
 \tag{F1T27.4}
\]
All entries $A,B,\beta$ are positive for $n\ge22$.
Here is a full verification of the bound.

For $(n-3,3)$, the cases $f=0,1,2$ give respectively
$-k_2+k_3$, $k_3$, and $-1+k_2+k_3$, each of absolute
value at most $B=n/2$. For $f\ge3$ the character is
at least $-2$. Using $2k_2+3k_3\le n-f$ gives the upper
bound
\[
 \chi(s)\le P(f):=(f-1)(f^2-8f+3n)/6.
\]
The function $P$ is convex on $[3,n-2]$, since
$P''(f)=f-3$. At the two endpoints,
\[
 P(3)=n-5,\quad P(n-2)=(n-3)A,\quad
 3A+B-P(3)=(n^2-10n+30)/2>0.
\]
The affine function $Af+B$ is above both endpoints,
and therefore above this convex function throughout the
interval. It also exceeds two, proving the absolute bound.

For $(n-3,1^3)$ the first three cases are
$k_2+k_3-1$, $k_3$, and $-k_2+k_3$, again bounded
in absolute value by $n/2$. For $f\ge3$ put
$Q(f)=\binom{f-1}3+(n-f)/3$. Then
\[
 -(f-1)n/2\le\chi(s)\le Q(f),\qquad Q''(f)=f-2.
\]
Since $A-n/2=(n-10)(n-2)/6\ge0$, the lower bound
is no smaller than $-(Af+B)$. The upper bound is convex,
and its two endpoint comparisons are
\[
 Q(3)=(n-3)/3\le3A+B,\qquad
 A(n-2)+B-Q(n-2)=(n^2-6n+16)/6>0.
\]
This proves the second affine estimate.

For $(n-3,2,1)$ the cases $f=0,1,2,3$ give
$-k_3$, $1-k_3$, $-k_3$, and $-1-k_3$.
Their absolute values are at most $n/3$.
For $4\le f\le n-2$, the polynomial
$f(f-2)(f-4)/3$ is nonnegative, and
$(f-2)(f-4)/3\le(n-4)(n-6)/3=A$.
Consequently $-n/3\le\chi(s)\le Af$, as required.

\paragraph{2. The actual cyclic projectors and every scalar phase.}
Let $g$ be a nonidentity original inertia image, put
$s=\sigma(g)$, $z=c(g)$, and let
$d=n-\#\{\hbox{cycles of }s\}$.
First suppose $s\ne1$, of order $q$.
If $z^q=1$, the actual projector onto its fixed vectors is
\[
  P_g=\frac1q\sum_{j=0}^{q-1}\rho(g)^j
      =\frac1q\sum_{j=0}^{q-1}z^jR_\lambda(s^j).
 \tag{F1T27.5}
\]
All nonzero powers $s^j$, $0<j<q$, are nonidentity.
The cyclic orbit-counting identity for the original
$n$ letters says
$q^{-1}\sum_{j=0}^{q-1}f(s^j)=n-d$.
Taking the trace of (F1T27.5), using (F1T27.4), and
then subtracting from $m$ gives
\[
\begin{split}
 \dim\ker(\rho(g)-1)&\le A(n-d)+B+\beta/q,\\
 a(g):=\operatorname{rank}(\rho(g)-1)
     &\ge Ad+(1-q^{-1})\beta
      \ge Ad+\beta/2.
\end{split}
 \tag{F1T27.6}
\]
This is the fixed subspace of the actual inertia; no
independence of $z$ and $s$ is assumed.

If $z^q\ne1$, any fixed vector would satisfy
$v=z^qv$, so $a(g)=m$. Since $d\le n-1$, the same
last inequality holds: the differences
$m-A(n-1)-\beta/2$, in the order of (F1T27.4), are
\[
 \frac{5(n-4)(n-2)}{12},\quad
 \frac{5n^2-24n+34}{12},\quad
 \frac{6n^2-35n+48}{6},
 \tag{F1T27.7}
\]
all positive for $n\ge22$.
Finally, if $s=1$, nonidentity of the original inertia
means $z\ne1$. Its moving rank is $m\ge\beta/2$ and
$d=0$. Thus (F1T27.6) holds for every original branch
entry, including the pure scalar ones.

The explicit sign-conjugate intertwiner of T22--T26
replaces $c$ by the actual character
$c\,\operatorname{sgn}\sigma$. It leaves all permutation
indices unchanged and is still a $\mu_4$ character.
Hence the same proof applies to every conjugate partition.

\paragraph{3. The original quotient curve and the strict contradiction.}
Let $C\to C/G=\mathbb P^1$ be the original cover,
and let $H_1=\sigma^{-1}(\operatorname{Stab}(1))$.
The actual connected quotient $X=C/H_1\to\mathbb P^1$
has degree $n$. For the original branch tuple
$g_1,\ldots,g_b$ its local ramification contribution is
exactly the permutation index $d_j$.
Riemann--Hurwitz therefore gives
\[
 \sum_jd_j=2n-2+2g(X)\ge2n-2.
 \tag{F1T27.8}
\]
Every branch is retained; pure scalar branches contribute
zero to this displayed sum. The original coefficient
Chevalley--Weil identity and period requirement are
$\sum_j a(g_j)=2m+6$ and $b\ge12$.
Summing (F1T27.6) now gives
\[
 \sum_j a(g_j)\ge2A(n-1)+6\beta>2m+6.
 \tag{F1T27.9}
\]
The three exact positive excesses are, respectively,
\[
 \frac{5n^2-30n-38}{3},\quad
 \frac{5n^2-18n-50}{3},\quad
 \frac{2(7n^2-25n-33)}3 .
 \tag{F1T27.10}
\]
Their values and derivatives are positive at $n=22$,
and their leading coefficients are positive. This proves
the uniform exclusion (F1T27.1).

\paragraph{4. An elementary dimension bound after T26.}
For completeness, let $D_3(n)=n(n-1)(n-5)/6$.
For every partition with at least three boxes outside
both its first row and its first column, and $n\ge15$,
\[
                        \dim S^\lambda\ge D_3(n).
 \tag{F1T27.11}
\]
Here is a direct tableau proof, without a minimal-degree
classification. The corner recurrence for standard
tableaux gives the dimension of the actual Specht
module, by the Frobenius and Young-rule proof already
read in T25. Exhaustive corner recurrence for $n=15,16$,
independently compared with the Frobenius dimension
formula, verifies the two induction bases.
The minima are $D_3(15)=350$ and $D_3(16)=440$;
all partitions, including self-conjugate ones, are checked.

On the boundary, a tail of length exactly three gives
one of the three shapes in (F1T27.3), or its conjugate.
Their dimensions exceed $D_3(n)$ by
$0$, $n-1$, and $n(n^2-6n+11)/6$, respectively.
In the interior both tails are at least four.
A nonrectangle has at least two removable corners,
and each child still has both tails at least three.
Their dimensions sum to at least $2D_3(n-1)$, and
\[
 2D_3(n-1)-D_3(n)=(n-1)(n-3)(n-8)/6>0.
\]
A rectangle with both side lengths at least two has
one child and two distinct grandchildren. Each side
length is at most $n/2$, so, for $n\ge17$, both tails
of each grandchild are at least $n/2-2\ge3$.
The two grandchildren therefore contribute at least
$2D_3(n-2)$. The remaining difference is
\[
 2D_3(n-2)-D_3(n)
       =(n^3-18n^2+77n-84)/6>0\quad(n\ge17);
\]
on writing $n=17+x$ its numerator is
$x^3+33x^2+332x+936$.
This completes the induction.
T26 removes all tails of size at most two in the
residual source, while T25 gives $n\ge22$.
Since $D_3$ is increasing there, (F1T27.2) follows.
The present stronger exclusion also removes the
boundary cases of this dimension argument.

\paragraph{Actual geometric and cohomological scope.}
The quotient used in (F1T27.8) is the literal
$C\to C/H_1$, with pullback and norm on Jacobians
satisfying $p_*p^*=[|H_1|]$ and
$p^*p_*=\sum_{h\in H_1}h$.
It supplies a ramification inequality; the sixfold
remains the original primitive-corner image in $J(C)$.
The only new coefficient endomorphism is the cyclic
projector (F1T27.5) on the original local system.
All relative-to-absolute cohomology maps, cup products,
Hermitian forms and primitive corner maps are those
of T24--T25, including the actual rank-one intertwiner
and its positive rational pullback scale.
Thus the same actual abelian family and cycle-class,
Hilbert/Chow, Poincare-biextension and Adams/GRR maps
are retained. No whole quotient Jacobian is substituted
for its coefficient factor, and no determinant formula
is inferred from the affine character estimate.

The primary character input is Guralnick--Tiep,
\emph{A problem of Kollar and Larsen on finite linear
groups}, JEMS 14 (2012), 605--657, Lemma 3.3 and its full
proof, printed pp.625--626, DOI \nolinkurl{10.4171/JEMS/313}.
The affine estimates, actual cyclic-projector bound
and source contradiction are proved above.
The accompanying certificate checks both dimension
induction bases and all permutation/scalar profiles
in two residual degrees; the uniform conclusion
comes from the proof, not from extrapolating those checks.
